\section{Liquidated damages mechanism}

\subsection{Definition and trigger}
Liquidated damages are a pre-agreed sum, stated in the supply contract, payable
by the OEM when a measured performance metric \emph{exceeds} its contractually
guaranteed threshold over the monitoring period. They provide commercial
certainty and avoid disputes over the quantum of actual loss.

Assessment is typically \textbf{per unit of excess}:
\begin{itemize}\setlength{\itemsep}{0.2em}
  \item \textbf{FOR} --- per additional forced outage above the guarantee.
  \item \textbf{FEU} --- per increment (e.g.\ each additional 0.1\,\%) of forced
        energy unavailability above the guarantee.
\end{itemize}

\subsection{Market ranges}
Indicative LD rates for a HVDC scheme of order 1\,GW:

\begin{center}
\renewcommand{\arraystretch}{1.35}
\begin{tabularx}{\linewidth}{@{}l l X@{}}
  \toprule
  \textbf{Metric} & \textbf{Penalty rate} & \textbf{Notes} \\
  \midrule
  FOR (per outage) & \$0.1--1.2\,M & per additional forced outage \\
  FEU              & \$75--600\,k  & per additional 0.1\,\% unavailability \\
  \textbf{LD cap}  & 3--5\,\%      & of contract value \\
  \bottomrule
\end{tabularx}
\end{center}

\subsection{Market context}
\begin{itemize}\setlength{\itemsep}{0.2em}
  \item The high end of each range reflects demanding customers in a buyer's
        market; the low end reflects supplier-favourable conditions.
  \item Caps have compressed over time: 3--5\,\% today versus 10--20\,\% in the
        buyer's market of $\sim$20 years ago.
  \item LDs are a \emph{sole-remedy} cap on operational performance in most
        contracts; the interaction with defect-liability and termination
        remedies must be checked case by case.
\end{itemize}
